\documentclass[10pt]{article}
\usepackage[utf8]{inputenc}
\usepackage[activeacute,spanish,es-nodecimaldot]{babel}
\usepackage[left=1.5cm,top=1.5cm,right=1.5cm, bottom=1.5cm,letterpaper, includeheadfoot]{geometry}
%\usepackage[parfill]{parskip}

\usepackage{amssymb, amsmath, amsthm}
\usepackage{graphicx}
\usepackage{lmodern,url}
\usepackage{paralist} %util para listas compactas

% paquetes agregados por mi
% para escribir algoritmos
\usepackage[spanish]{algorithm}
\usepackage[noend]{algpseudocode}


\usepackage{tikz}
%\usetikzlibrary{positioning,chains,fit,shapes,calc}
\usetikzlibrary{positioning, shapes, arrows,decorations.text}
\usepackage{float}
\usepackage{caption}
\usepackage{ mathdots }
\usepackage{verbatim}
\usepackage{ stmaryrd }
\usepackage{fancyhdr}
\pagestyle{fancy}
\fancypagestyle{plain}{%
\fancyhf{}
\lhead{\footnotesize\itshape\bfseries\rightmark}
\rhead{\footnotesize\itshape\bfseries\leftmark}
}


% macros
\newcommand{\Q}{\mathbb Q}
\newcommand{\R}{\mathbb R}
\newcommand{\N}{\mathbb N}
\newcommand{\Z}{\mathbb Z}
\newcommand{\C}{\mathbb C}

% macros jp
\newcommand{\cev}[1]{\reflectbox{\ensuremath{\vec{\reflectbox{\ensuremath{#1}}}}}}
\newcommand{\sfunction}[1]{\textsf{\textsc{#1}}}
\algrenewcommand\algorithmicforall{\textbf{foreach}}
\algrenewcommand\algorithmicindent{.8em}
\newcommand{\val}[1]{\text{val}(#1)}

%Teoremas, Lemas, etc.
\theoremstyle{plain}
\newtheorem{teo}{Teorema}
\newtheorem{lem}{Lema}
\newtheorem{prop}{Proposición}
\newtheorem{cor}{Corolario}

\theoremstyle{definition}
\newtheorem{defi}{Definición}
% fin macros

%%%%% NOMBRE ESCRIBAS Y FECHA
\newcommand{\sca}{Escriba Uno}
\newcommand{\scb}{Escriba Dos}
\newcommand{\scc}{Escriba Tres}
\newcommand{\catnum}{0} %numero de catedra
\newcommand{\fecha}{5 de noviembre 2018 }

%%%%%%%%%%%%%%%%%%

%Macros para este documento
\newcommand{\cin}{\operatorname{cint}}

\begin{document}
%Encabezado
\fancyhead[L]{Facultad de Ciencias Físicas y Matemáticas}
\fancyhead[R]{Universidad de Chile}
\vspace*{-1.2 cm}
%\begin{minipage}{0.6\textwidth}
%\begin{flushleft}
%\hspace*{-0.5cm}\textbf{MA4702. Programación Lineal Mixta. 2018.}\\
%\hspace*{-0.5cm}\textbf{Profesor:} José Soto\\
%\hspace*{-0.5cm}\textbf{Escriba(s):} \sca, \scb~y \scc.\\
%\hspace*{-0.5cm}\textbf{Fecha:} \fecha.
%\end{flushleft}
%\end{minipage}
\begin{flushleft}
\includegraphics[scale=0.15]{fcfm}
\end{flushleft}
\bigskip
%Fin encabezado


\subsection{Algunos resultados de aproximación para ATMP (caso ganancias lineales)}

Sea $P$ el problema ATMP (asginaci\'on de trabajos en m\'aquinas paralelas). Se estudiar\'an algoritmos de aproximaci\'on para el caso en que las curvas de ganancia para cada trabajo $i \in [n]$ tienen la forma:


%\begin{center}
%\textit{Se tienen $n$ trabajos y $m$ m\'aquinas. Cada trabajo $i \in [n]$ tiene asociado una tasa de ganancia $p_i$ y un tiempo $x_i$ tales que la ganancia $g_i(t)$ obtenida por procesar dicho trabajo $t$ unidades de tiempo est\'a dada por}
$$
g_i(t) = \left\{\begin{array}{lr}
p_i t & t \leq t_i \\
p_i t_i & t>t_i
\end{array}\right.
$$
donde $p_i$ es la tasa de ganancia asociada al trabajo $i-$\'esimo, y $t_i$ es el tiempo de proceso m\'aximo despu\'es del cual el trabajo $i$ deja de aportar beneficio. Suponagmos que los trabajos est\'an indexados de modo que:
\begin{equation}
\label{eq:trabajosOrd}
p_1 \geq p_2 \geq \cdots \geq p_n
\end{equation}

%\textit{El objetivo es encontrar una asignaci\'on de trabajos a m\'aquinas de m\'aximo beneficio tal que:}
%\begin{enumerate}
%\item Sea $t_j$ el instante en que termina el \'ultimo trabajo la m\'aquina $j-$\'esima. Entonces $t_j \leq 1$.
%
%\item 
%\end{enumerate}
%\end{center}



%Sea $\val{\text{OPT}}$ el valor de alguna asignación \'optima para el problema $P$. A continuación se enuncian un par de proposiciones que serán \'utiles m\'as adelante (la demostración del segundo queda pendiente).

\begin{defi}[Asignación]
Una asignación $\mathcal{A}$ es una tupla $(A,T,M,X)$ donde:
\begin{itemize}
\item $A \subseteq [n]$ es el conjunto de \textit{trabajos asignados}.
\item $T = \{y_i\}_{i \in A}$ son los \textit{tiempos de proceso} de cada trabajo $i \in A$.
\item $M = \{m_i\}_{i \in A}$ es el conjunto de máquinas asignadas para cada trabajo en $A$.
%\item $X$ es el conjunto de \textit{posiciones} de procesamiento de cada trabajo $i \in A$ en $m_i$, i.e. $i$ es el $x_i-$ésimo trabajo procesado en la máquina $m_i$.
\end{itemize}
\end{defi}

Se denotará por $\val{\mathcal{A}}$ al \textit{valor} de la asignación $\mathcal{A}$, y diremos que un trabajo fue \textit{cortado} según la asignación $\mathcal{A}$ si $0<y_i < t_i$. 
\begin{lem}
\label{lem:cortes}
Sea $\mathcal{A}$ una asignación factible. Entonces se puede obtener una asignación $\mathcal{A}'$ a partir de $\mathcal{A}$ de modo que solo se haya cortado un trabajo por máquina, y $\val{\mathcal{A}'} \geq \val{\mathcal{A}}$.
\end{lem}
\begin{proof}
Sean $i,j$ dos trabajos cortados según $\mathcal{A}$, sin pérdida de generalidad supongase que  $p_i \geq p_j$. El tiempo disponible que le queda al trabajo $i$ para poder ser procesado y aumentar su contribución al beneficio está dado por $t_i-y_i$. Luego, basta hacer los cambios $y_i' \gets y_i + \min\{t_i-y_i,y_j\}$, $y_j' \gets y_j -\min\{t_i-y_i,y_j\}$, y definir $T'$ con estos nuevos tiempos de proceso $y_i'$ y $y_j'$ para los trabajos $i$ y $j$, respectivamente. Definiendo $\mathcal{A}'$ con estos nuevos tiempos de proceso, se tiene que, o bien $i$ se realiza en $t_i$, o bien el trabajo $j$ se deja de realizar. Luego el numero de trabajos cortados disminuye al menos en $1$. Por otro lado, dado que $p_i \geq p_j$ se tiene que:
$$
\val{\mathcal{A}'} \geq \val{\mathcal{A}}
$$
\end{proof}

Otro resultado de utilidad es el siguiente.
\begin{lem}
\label{lem:caracOPT}
Sea $\mathcal{A} = (A,T,M)$ alguna asignación óptima para el problema $P$. Se tiene entonces que: 
\begin{equation}
\label{eq:caractOPT}
p_i \geq p_j \quad \forall i \in A,\ j \in [n]\setminus{A}
\end{equation}
i.e., el óptimo asigna los trabajos con mejor pendiente.
\end{lem}
\begin{proof}
Por contradicción, supongase que existen $i\in A,j \in [n]\setminus{A}$ tales que no se cumple (\ref{eq:caractOPT}). Entonces haciendo el cambio $y_i \gets 0$ y $y_j \gets \min\{y_i,t_j\}$, se tendría un asignamiento $\mathcal{A}'$ tal que:
$$
\val{\mathcal{A}'} > \val{\mathcal{A}}
$$
lo cual contradide la optimalidad de  $\mathcal{A}$.
\end{proof}
De los lemas (\ref{lem:cortes}) y (\ref{lem:caracOPT}) se concluye directamente el siguiente resultado:
\begin{prop}
\label{prop:final}
Existe una asignación $\mathcal{A} = (A,T,M)$ óptima para $P$ que cumple las siguientes características:
\begin{enumerate}
\item  $p_i \geq p_j$, para $i \in A$, $j \in [n]\setminus{A}$.
\item Toda máquina corta a lo más un trabajo.
\end{enumerate}
\end{prop}

\textbf{Obs:} De la proposición (\ref{prop:final}) y de el orden de las pendientes en (\ref{eq:trabajosOrd}) se puede deducir que existe una asignación óptima tal que el conjunto de trabajos asignados tiene la forma $A = [\ell]$ para algún $\ell \in [n]$.\\~\\

%\begin{proof}
%En efecto, considere una máquina $j$ y supongamos que existen dos trabajos $j_1$ y $j_2$ tales que cada uno se procesa un tiempo $h_1 < x_{j_1}$ y $h_2 < x_{j_2}$, respectivamente. Sin perdida de generalidad, supongamos que $p_1 \geq p_2$. Llamese por $A$ la presente asignaci\'on. Se separar\'a el an\'alisis en diferentes casos:
%\begin{itemize}
%\item \textbf{caso $h_2 \geq x_{j_1}-h_1$}:
%
%En este caso, como $p_1 \geq p_2$, basta con reemplazar el tiempo de proceso de $j_2$ por tiempo de proceso de $j_1$, de modo que $j_1$ no sea cortado, y $j_2$
%\end{itemize}
%\end{proof}

Para encontrar un algoritmo de aproximaci\'on, dividiremos el problema según el \textit{largo} de los trabajos. Para ello, sea $k \in \N$ fijo.
\subsubsection*{caso $x_i \leq 1/k \quad \forall i \in [n]$}

%Sea $P$ el problema ATMP, y considerese $P'$ el problema donde, en vez de tener la restricción de dead line igual a $1$, ahora solo se pide que cada trabajo asignado haya partido antes de $1$. Además, considere que los largos de cada trabajo cumplen la relación $x_i \leq 1/k$, para algún $k \in \N$ fijo. Como antes, hay $n$ trabajos a agendar y $m$ el número de máquinas. Pir último, considere el siguiente algoritmo:

%Para el problema $P$, sea $\text{OPT}$ el valor de la asignación optima. A continuación enunciamos un par de proposisiones que serán utiles:

Para el presente caso se mostrar\'a que un algoritmo glot\'on logra una $(1+1/k)^2$ aproximaci\'on. Considere para ello una variante del problema $P$, llamada $P'$, donde se realizan los siguientes cambios en las restricciones: 
\begin{enumerate}
\item Todo asignación $\mathcal{A}$ tiene que conjunto de trabajos asignados de la forma $A=[\ell]$, para algún $\ell \in [n]$. 
\item Para todo $i \in A$, $y_i = t_i$.
\item Todo trabajo agendado debe partir antes de $1$. 
\end{enumerate}

Para describir el algoritmo Considere las siguientes notaciones. Sea $j \in [m]$ una m\'aquina cualquiera:
\begin{itemize}
\item $L_j$: instante en que la máquina $j$ termina el \'ultimo trabajo.
%\item $t$: instante m\'as chico en que se termina el \'ultimo trabajo asignado.
\item $\ell$: \'ultimo trabajo asignado.
\item $m_\ell$: m\'aquina a la que se asigna el trabajo $\ell$ (un valor de $0$ indica que no ha sido asignado).
\item $x_\ell$: instante en el que el trabajo $\ell$ parte.
\end{itemize}
Con esto se puede describir un algoritmo glot\'on para $P'$:
\begin{algorithm}[H]
\caption{Glotón para $P'$}\label{alg:greddyP'}
\begin{algorithmic}[1]

\State ordenar las pendientes de modo que $p_1 \geq p_2 \geq p_3 \geq \cdots \geq p_n$  \;
\State $\ell \gets 0$\;
\For{$i: 1 \to n$}
	\State $L_i \gets 0$\;
	\State $m_i \gets 0$\;
\EndFor
\While {$\min\{L_j:j \in [m]\}\leq 1$}
	\State $\ell \gets \ell+1$\;
	\State $j \gets \text{argmin}\{L_j: j \in [m]\}$\;
	\State $m_\ell \gets j$\;
	\State $L_j \gets \sum_{i = 1}^n \llbracket m_i=m_\ell\rrbracket t_i + t_\ell$
	\State $x_\ell \gets L_j-t_\ell$\;
\EndWhile
\State \textbf{end}
\State \textbf{Return } $(A = [\ell], X = \{x_i\}_{i=1}^\ell, M = \{m_i\}_{i=1}^{\ell})$
\end{algorithmic}
\end{algorithm}


Se denotará por ALG al algoritmo (\ref{alg:greddyP'}). Antes de probar que ALG induce una $(1+1/k)^2$ aproximación para el problema $P$, un resultado que es de utilidad:

\begin{prop}
\label{prop:appP'}
ALG es una $1+1/k$ aproxeimación para el problema $P'$.
\end{prop}
\begin{proof}
En efecto, sea $\mathcal{A}_{P'}$ una asignación óptima para $P'$, y sea $\mathcal{A}_{\text{ALG}}$ la asignación entregada por ALG, y considere $A_{P'} = [\ell_{P'}]$ y $A_{\text{ALG}}= [\ell_{\text{ALG}}]$. Dada la condición (\ref{eq:trabajosOrd}), sabemos el último trabajo asignado por ALG cumple que
\begin{equation}
\label{eq:minorante}
p_{\ell_{\text{ALG}}} \leq p_i \quad \forall i \in [\ell_{\text{ALG}}]
\end{equation}
Por otro lado, es claro que $\ell_{\text{ALG}} \leq \ell_{\text{OPT}}$. Denotemos por $\val{\text{ALG}}$ y $\val{\text{OPT}}$ los valores de la función objetivo para los agendamientos de ALG y OPT respectivamente. Entonces:
\begin{align*}
\val{\text{OPT}} &= \sum_{i = 1}^{\ell_{\text{OPT}}} p_i t_i \\
		  &= \sum_{i=1}^{\ell_\text{ALG}}p_i t_i + \sum_{i = \ell_{\text{ALG}}+1}^{\ell_{\text{OPT}}}p_it_i\\
		  &\leq \val{\text{ALG}}+p_{\ell_\text{ALG}} \sum_{i = \ell_{\text{ALG}}+1}^{\ell_{\text{OPT}}}t_i \\
		  &\leq \val{\text{ALG}}+p_{\ell_\text{ALG}} + p_{\ell_\text{ALG}} m\frac{1}{k} \qquad (\star)
\end{align*}
donde $(\star)$ por otro lado, notemos que:
\begin{align*}
\val{\text{ALG}} &= \sum_{i=1}^{\ell_{\text{ALG}}} p_i t_i \\
		  &\geq p_{\ell_{\text{ALG}}}\sum_{i=1}^{\ell_{\text{ALG}}} t_i\\
		  &\geq p_{\ell_{\text{ALG}}}m
\end{align*}
Así, juntando los dos desarrollos anteriores se concluye que:
$$
\val{\text{OPT}} \leq \left(1 + \frac{1}{k}\right) \val{\text{ALG}}
$$
\end{proof}
De la proposcición (\ref{prop:appP'}) se deduce el siguiente corolario:
\begin{cor}
ALG induce una asignación para el problema $P$ que es una $(1+1/k)^2$ aproximación.
\end{cor}
\begin{proof}
Sean $\val{\text{OPT}^P}$ y $\val{\text{OPT}^{P'}}$ los valores óptimos de $P$ y $P'$, respectivamente. Más aún, consideremos la asignación óptima de $P$ con la estructura descrita en la observación de la proposición (\ref{prop:final}). Dicha asignación para $P$ tiene beneficio dado por $\val{\text{OPT}^P}$, y ademas cumple con ser una asignación factible para $P'$, luego:
\begin{equation}
\label{eq:PP'}
\val{\text{OPT}^P} \leq \val{\text{OPT}^{P'}}
\end{equation}

gracias a la proposición (\ref{prop:appP'}), se tiene además que:
$$
\val{\text{OPT}^P} \leq \val{\text{OPT}^{P'}} \leq \left(1+\frac{1}{k} \right)\val{\text{ALG}}
$$

Por otro lado, el beneficio $\val{\text{ALG}}$ puede ser dividido en el valor obtenido por los trabajos procesados hasta el instante $1$, y el valor asociado a los que se procesaron en instantes posteriores
\begin{equation}
\label{eq:divP'}
\val{\text{ALG}} = \val{\text{ALG}_{>1}} + \val{\text{ALG}_{\leq 1}}
\end{equation}

Se demostrar\'a entonces que
\begin{equation}
\label{eq:cotaP'}
\val{\text{ALG}_{>1}} \leq \frac{1}{k}\val{\text{ALG}_{\leq 1}}
\end{equation}

Notar que de (\ref{eq:cotaP'}), (\ref{eq:divP'}) y (\ref{eq:PP'}) se deduce directamente que:
\begin{equation*}
\val{\text{OPT}^P} \leq \left( 1 + \frac{1}{k} \right)^2\val{\text{ALG}_{\leq 1}}
\end{equation*}

Así, basta probar (\ref{eq:cotaP'}) para concluir el resultado. Para ello, considere por un momento los trabajos asignados a una máquina $j \in [m]$, y abusando de notación llamemos $p_1 \ldots p_l$ las pendientes de los trabajos asignados a dicha máquina, realizados en el mismo orden de indexación. El tiempo que ALG procesa el último trabajo en la máquina $j$, $t_\ell$, puede ser escrito como $t_\ell = \tilde{t}_\ell + t'_\ell$, donde $\tilde{t}_\ell$ es la cantidad de tiempo que ALG procesa $\ell$ hasta el instante $1$, y $t'_\ell$ es el tiempo que se procesa $\ell$ después del instante $1$. Sea $\val{\text{ALG}_{j,\leq 1}}$ el valor asociado a los trabajos que se procesan hasta el instante $1$ en la máquina $j$, i.e.
\begin{equation}
\label{eq:ganjP}
\val{\text{ALG}_{j,\leq 1}}  = \sum_{i <\ell} p_i t_i + p_\ell \tilde{t}_\ell
\end{equation} 
Luego, notemos que $1 = \sum_{i<\ell} t_i + \tilde{t}_{\ell}$, de donde

\begin{equation*}
 p_\ell(1-\tilde{t}_\ell) = p_\ell\sum_{i < \ell} t_i = \sum_{i < \ell} p_\ell t_i \leq  \sum_{i < \ell} p_i t_i
\end{equation*}
donde la última desigualdad viene del hecho de que ALG asigna los trabajos en orden decreciente de pendiente, de modo que $p_\ell \leq p_i \ \forall i \in [\ell]$. Del último desarrollo se tiene que
\begin{equation}
\label{eq:cotasGan}
p_\ell \leq \sum_{j < \ell} p_j t_j + p_\ell \tilde{t}_\ell
\end{equation}
usando que $t_\ell \leq 1/k$, se deduce que
\begin{equation}
\label{eq:cotaAntesDps}
p_\ell t'_\ell \leq \frac{1}{k}\left(\sum_{j < \ell} p_j t_j + p_\ell \tilde{t}_\ell \right) = \frac{1}{k}\val{\text{ALG}_{j,\leq 1}}
\end{equation}
Así, considere ahora $\ell_1, \ldots, \ell_m$ los \'ultimos trabajos que cada m\'aquina procesa seg\'un ALG, y sus respectivos $\{p_{\ell_i}\}_{i=1}^m$ y $\{t_{\ell_i}\}_{i=1}^m$. Notar que
$$
\val{\text{ALG}_{>1}} = \sum_{i = 1}^m p_{\ell_i}t'_{\ell_i}, \quad   \val{\text{OPT}_{\leq 1}}  = \sum_{i=1}^m \text{ALG}^{P'}_{j_i,\leq 1}  
$$
Entonces, por la desigualdad (\ref{eq:cotaAntesDps}) se puede deducir que
\begin{equation*}
\val{\text{ALG}_{>1}} = \sum_{i = 1}^m p_{\ell_i} t'_{\ell_i} \leq \sum_{i=1}^m \frac{1}{k} \val{\text{ALG}_{j_i,\leq 1}} = \frac{1}{k} \val{\text{ALG}_{\leq 1}}
\end{equation*}
lo que termina la demostración.
\end{proof}
\subsubsection*{caso $x_i >1/k$}
\newpage

%Para este caso, se puede resolver el problema de manera exacta en tiempo polinomial si se consideran $m$ y $k$ constantes. Para ello basta notar que la cantidad de trabajos que puede tener cada m\'aquina es menor que $km$, y entonces la cantidad de asignaciones factibles para $P$ en este caso es $O(m^{km})$
\end{document}